% Estilo de listagens de código (pacote listings), usado pelos exemplos desta pasta.
% Não é um documento: é incluído com % Estilo de listagens de código (pacote listings), usado pelos exemplos desta pasta.
% Não é um documento: é incluído com % Estilo de listagens de código (pacote listings), usado pelos exemplos desta pasta.
% Não é um documento: é incluído com % Estilo de listagens de código (pacote listings), usado pelos exemplos desta pasta.
% Não é um documento: é incluído com \input{estilo-codigo} no preâmbulo.
% Testado com: listings, TeX Live 2024 ou mais recente. O trecho marcado com o id "estilo" é o que as receitas e o livro exibem.
%%BEGIN:estilo
\usepackage[varqu,varl,scaled=0.95]{inconsolata} % largura fixa legível (0/O, 1/l/I e & distintos)
\usepackage{xcolor}
\usepackage{listings}
\lstloadaspects{tex,emph}               % habilita as chaves texcsstyle, moretexcs e emphstyle
\definecolor{cod-kw}{RGB}{0,51,153}     % palavras-chave e comandos
\definecolor{cod-com}{RGB}{0,110,50}    % comentários
\definecolor{cod-str}{RGB}{163,21,21}   % cadeias de caracteres
\definecolor{cod-foco}{RGB}{190,70,0}   % palavras em destaque (opção emph)
\lstdefinestyle{codigo}{
  basicstyle=\ttfamily\footnotesize, columns=fixed, basewidth=0.5em,
  keepspaces=true, tabsize=4,
  keywordstyle=\bfseries\color{cod-kw}, commentstyle=\itshape\color{cod-com},
  stringstyle=\color{cod-str}, showstringspaces=false,
  texcsstyle=*\bfseries\color{cod-kw}, emphstyle=\bfseries\color{cod-foco},
  % comandos de pacotes que a linguagem [LaTeX]TeX não conhece
  moretexcs={includegraphics, caption, toprule, midrule, bottomrule, cmidrule,
    multicolumn, multirow, IBGEtab, fonte, nota, legenda, autocite, textcite,
    cite, parencite, printbibliography, addbibresource, autoref, cref, Cref,
    eqref, ref, pageref, url, href, lstinputlisting, lstset, lstdefinestyle,
    lstinline, frametitle, pause, only, uncover, includeonly, input, include,
    maketitle, tableofcontents, listoffigures, listoftables, newcommand,
    renewcommand, setlength, textwidth, linewidth, hline, cline, footnote,
    textit, texttt, frac, sqrt, autor, titulo, orientador, instituicao,
    local, data, textual, postextual, pretextual, imprimircapa,
    imprimirfolhaderosto, acrodef, ac, gls, question, part, choice,
    CorrectChoice, points},
  numbers=left, numberstyle=\tiny\color{black!55}, numbersep=7pt,
  backgroundcolor=\color{black!4}, frame=single, rulecolor=\color{black!25},
  breaklines=true, postbreak={\mbox{\textcolor{cod-com}{$\hookrightarrow$}\space}},
  extendedchars=true,
  literate={á}{{\'a}}1 {é}{{\'e}}1 {í}{{\'i}}1 {ó}{{\'o}}1 {ú}{{\'u}}1
           {Á}{{\'A}}1 {É}{{\'E}}1 {Í}{{\'I}}1 {Ó}{{\'O}}1 {Ú}{{\'U}}1
           {â}{{\^a}}1 {ê}{{\^e}}1 {ô}{{\^o}}1 {ã}{{\~a}}1 {õ}{{\~o}}1
           {Â}{{\^A}}1 {Ê}{{\^E}}1 {Ô}{{\^O}}1 {Ã}{{\~A}}1 {Õ}{{\~O}}1
           {ç}{{\c c}}1 {Ç}{{\c C}}1 {à}{{\`a}}1 {À}{{\`A}}1
           {ü}{{\"u}}1 {Ü}{{\"U}}1
           {º}{{\textordmasculine}}1 {ª}{{\textordfeminine}}1,
}
% versão em preto e branco do mesmo estilo (para impressão sem cores)
\lstdefinestyle{codigopb}{style=codigo, keywordstyle=\bfseries,
  texcsstyle=*\bfseries, emphstyle=\bfseries\itshape,
  commentstyle=\itshape, stringstyle=\relax, backgroundcolor=\color{white},
  rulecolor=\color{black}, numberstyle=\tiny}
\lstset{style=codigo}
\renewcommand{\lstlistingname}{Código}            % rótulo da legenda: "Código N"
\renewcommand{\lstlistlistingname}{Lista de códigos}
%%END:estilo
 no preâmbulo.
% Testado com: listings, TeX Live 2024 ou mais recente. O trecho marcado com o id "estilo" é o que as receitas e o livro exibem.
%%BEGIN:estilo
\usepackage[varqu,varl,scaled=0.95]{inconsolata} % largura fixa legível (0/O, 1/l/I e & distintos)
\usepackage{xcolor}
\usepackage{listings}
\lstloadaspects{tex,emph}               % habilita as chaves texcsstyle, moretexcs e emphstyle
\definecolor{cod-kw}{RGB}{0,51,153}     % palavras-chave e comandos
\definecolor{cod-com}{RGB}{0,110,50}    % comentários
\definecolor{cod-str}{RGB}{163,21,21}   % cadeias de caracteres
\definecolor{cod-foco}{RGB}{190,70,0}   % palavras em destaque (opção emph)
\lstdefinestyle{codigo}{
  basicstyle=\ttfamily\footnotesize, columns=fixed, basewidth=0.5em,
  keepspaces=true, tabsize=4,
  keywordstyle=\bfseries\color{cod-kw}, commentstyle=\itshape\color{cod-com},
  stringstyle=\color{cod-str}, showstringspaces=false,
  texcsstyle=*\bfseries\color{cod-kw}, emphstyle=\bfseries\color{cod-foco},
  % comandos de pacotes que a linguagem [LaTeX]TeX não conhece
  moretexcs={includegraphics, caption, toprule, midrule, bottomrule, cmidrule,
    multicolumn, multirow, IBGEtab, fonte, nota, legenda, autocite, textcite,
    cite, parencite, printbibliography, addbibresource, autoref, cref, Cref,
    eqref, ref, pageref, url, href, lstinputlisting, lstset, lstdefinestyle,
    lstinline, frametitle, pause, only, uncover, includeonly, input, include,
    maketitle, tableofcontents, listoffigures, listoftables, newcommand,
    renewcommand, setlength, textwidth, linewidth, hline, cline, footnote,
    textit, texttt, frac, sqrt, autor, titulo, orientador, instituicao,
    local, data, textual, postextual, pretextual, imprimircapa,
    imprimirfolhaderosto, acrodef, ac, gls, question, part, choice,
    CorrectChoice, points},
  numbers=left, numberstyle=\tiny\color{black!55}, numbersep=7pt,
  backgroundcolor=\color{black!4}, frame=single, rulecolor=\color{black!25},
  breaklines=true, postbreak={\mbox{\textcolor{cod-com}{$\hookrightarrow$}\space}},
  extendedchars=true,
  literate={á}{{\'a}}1 {é}{{\'e}}1 {í}{{\'i}}1 {ó}{{\'o}}1 {ú}{{\'u}}1
           {Á}{{\'A}}1 {É}{{\'E}}1 {Í}{{\'I}}1 {Ó}{{\'O}}1 {Ú}{{\'U}}1
           {â}{{\^a}}1 {ê}{{\^e}}1 {ô}{{\^o}}1 {ã}{{\~a}}1 {õ}{{\~o}}1
           {Â}{{\^A}}1 {Ê}{{\^E}}1 {Ô}{{\^O}}1 {Ã}{{\~A}}1 {Õ}{{\~O}}1
           {ç}{{\c c}}1 {Ç}{{\c C}}1 {à}{{\`a}}1 {À}{{\`A}}1
           {ü}{{\"u}}1 {Ü}{{\"U}}1
           {º}{{\textordmasculine}}1 {ª}{{\textordfeminine}}1,
}
% versão em preto e branco do mesmo estilo (para impressão sem cores)
\lstdefinestyle{codigopb}{style=codigo, keywordstyle=\bfseries,
  texcsstyle=*\bfseries, emphstyle=\bfseries\itshape,
  commentstyle=\itshape, stringstyle=\relax, backgroundcolor=\color{white},
  rulecolor=\color{black}, numberstyle=\tiny}
\lstset{style=codigo}
\renewcommand{\lstlistingname}{Código}            % rótulo da legenda: "Código N"
\renewcommand{\lstlistlistingname}{Lista de códigos}
%%END:estilo
 no preâmbulo.
% Testado com: listings, TeX Live 2024 ou mais recente. O trecho marcado com o id "estilo" é o que as receitas e o livro exibem.
%%BEGIN:estilo
\usepackage[varqu,varl,scaled=0.95]{inconsolata} % largura fixa legível (0/O, 1/l/I e & distintos)
\usepackage{xcolor}
\usepackage{listings}
\lstloadaspects{tex,emph}               % habilita as chaves texcsstyle, moretexcs e emphstyle
\definecolor{cod-kw}{RGB}{0,51,153}     % palavras-chave e comandos
\definecolor{cod-com}{RGB}{0,110,50}    % comentários
\definecolor{cod-str}{RGB}{163,21,21}   % cadeias de caracteres
\definecolor{cod-foco}{RGB}{190,70,0}   % palavras em destaque (opção emph)
\lstdefinestyle{codigo}{
  basicstyle=\ttfamily\footnotesize, columns=fixed, basewidth=0.5em,
  keepspaces=true, tabsize=4,
  keywordstyle=\bfseries\color{cod-kw}, commentstyle=\itshape\color{cod-com},
  stringstyle=\color{cod-str}, showstringspaces=false,
  texcsstyle=*\bfseries\color{cod-kw}, emphstyle=\bfseries\color{cod-foco},
  % comandos de pacotes que a linguagem [LaTeX]TeX não conhece
  moretexcs={includegraphics, caption, toprule, midrule, bottomrule, cmidrule,
    multicolumn, multirow, IBGEtab, fonte, nota, legenda, autocite, textcite,
    cite, parencite, printbibliography, addbibresource, autoref, cref, Cref,
    eqref, ref, pageref, url, href, lstinputlisting, lstset, lstdefinestyle,
    lstinline, frametitle, pause, only, uncover, includeonly, input, include,
    maketitle, tableofcontents, listoffigures, listoftables, newcommand,
    renewcommand, setlength, textwidth, linewidth, hline, cline, footnote,
    textit, texttt, frac, sqrt, autor, titulo, orientador, instituicao,
    local, data, textual, postextual, pretextual, imprimircapa,
    imprimirfolhaderosto, acrodef, ac, gls, question, part, choice,
    CorrectChoice, points},
  numbers=left, numberstyle=\tiny\color{black!55}, numbersep=7pt,
  backgroundcolor=\color{black!4}, frame=single, rulecolor=\color{black!25},
  breaklines=true, postbreak={\mbox{\textcolor{cod-com}{$\hookrightarrow$}\space}},
  extendedchars=true,
  literate={á}{{\'a}}1 {é}{{\'e}}1 {í}{{\'i}}1 {ó}{{\'o}}1 {ú}{{\'u}}1
           {Á}{{\'A}}1 {É}{{\'E}}1 {Í}{{\'I}}1 {Ó}{{\'O}}1 {Ú}{{\'U}}1
           {â}{{\^a}}1 {ê}{{\^e}}1 {ô}{{\^o}}1 {ã}{{\~a}}1 {õ}{{\~o}}1
           {Â}{{\^A}}1 {Ê}{{\^E}}1 {Ô}{{\^O}}1 {Ã}{{\~A}}1 {Õ}{{\~O}}1
           {ç}{{\c c}}1 {Ç}{{\c C}}1 {à}{{\`a}}1 {À}{{\`A}}1
           {ü}{{\"u}}1 {Ü}{{\"U}}1
           {º}{{\textordmasculine}}1 {ª}{{\textordfeminine}}1,
}
% versão em preto e branco do mesmo estilo (para impressão sem cores)
\lstdefinestyle{codigopb}{style=codigo, keywordstyle=\bfseries,
  texcsstyle=*\bfseries, emphstyle=\bfseries\itshape,
  commentstyle=\itshape, stringstyle=\relax, backgroundcolor=\color{white},
  rulecolor=\color{black}, numberstyle=\tiny}
\lstset{style=codigo}
\renewcommand{\lstlistingname}{Código}            % rótulo da legenda: "Código N"
\renewcommand{\lstlistlistingname}{Lista de códigos}
%%END:estilo
 no preâmbulo.
% Testado com: listings, TeX Live 2024 ou mais recente. O trecho marcado com o id "estilo" é o que as receitas e o livro exibem.
%%BEGIN:estilo
\usepackage[varqu,varl,scaled=0.95]{inconsolata} % largura fixa legível (0/O, 1/l/I e & distintos)
\usepackage{xcolor}
\usepackage{listings}
\lstloadaspects{tex,emph}               % habilita as chaves texcsstyle, moretexcs e emphstyle
\definecolor{cod-kw}{RGB}{0,51,153}     % palavras-chave e comandos
\definecolor{cod-com}{RGB}{0,110,50}    % comentários
\definecolor{cod-str}{RGB}{163,21,21}   % cadeias de caracteres
\definecolor{cod-foco}{RGB}{190,70,0}   % palavras em destaque (opção emph)
\lstdefinestyle{codigo}{
  basicstyle=\ttfamily\footnotesize, columns=fixed, basewidth=0.5em,
  keepspaces=true, tabsize=4,
  keywordstyle=\bfseries\color{cod-kw}, commentstyle=\itshape\color{cod-com},
  stringstyle=\color{cod-str}, showstringspaces=false,
  texcsstyle=*\bfseries\color{cod-kw}, emphstyle=\bfseries\color{cod-foco},
  % comandos de pacotes que a linguagem [LaTeX]TeX não conhece
  moretexcs={includegraphics, caption, toprule, midrule, bottomrule, cmidrule,
    multicolumn, multirow, IBGEtab, fonte, nota, legenda, autocite, textcite,
    cite, parencite, printbibliography, addbibresource, autoref, cref, Cref,
    eqref, ref, pageref, url, href, lstinputlisting, lstset, lstdefinestyle,
    lstinline, frametitle, pause, only, uncover, includeonly, input, include,
    maketitle, tableofcontents, listoffigures, listoftables, newcommand,
    renewcommand, setlength, textwidth, linewidth, hline, cline, footnote,
    textit, texttt, frac, sqrt, autor, titulo, orientador, instituicao,
    local, data, textual, postextual, pretextual, imprimircapa,
    imprimirfolhaderosto, acrodef, ac, gls, question, part, choice,
    CorrectChoice, points},
  numbers=left, numberstyle=\tiny\color{black!55}, numbersep=7pt,
  backgroundcolor=\color{black!4}, frame=single, rulecolor=\color{black!25},
  breaklines=true, postbreak={\mbox{\textcolor{cod-com}{$\hookrightarrow$}\space}},
  extendedchars=true,
  literate={á}{{\'a}}1 {é}{{\'e}}1 {í}{{\'i}}1 {ó}{{\'o}}1 {ú}{{\'u}}1
           {Á}{{\'A}}1 {É}{{\'E}}1 {Í}{{\'I}}1 {Ó}{{\'O}}1 {Ú}{{\'U}}1
           {â}{{\^a}}1 {ê}{{\^e}}1 {ô}{{\^o}}1 {ã}{{\~a}}1 {õ}{{\~o}}1
           {Â}{{\^A}}1 {Ê}{{\^E}}1 {Ô}{{\^O}}1 {Ã}{{\~A}}1 {Õ}{{\~O}}1
           {ç}{{\c c}}1 {Ç}{{\c C}}1 {à}{{\`a}}1 {À}{{\`A}}1
           {ü}{{\"u}}1 {Ü}{{\"U}}1
           {º}{{\textordmasculine}}1 {ª}{{\textordfeminine}}1,
}
% versão em preto e branco do mesmo estilo (para impressão sem cores)
\lstdefinestyle{codigopb}{style=codigo, keywordstyle=\bfseries,
  texcsstyle=*\bfseries, emphstyle=\bfseries\itshape,
  commentstyle=\itshape, stringstyle=\relax, backgroundcolor=\color{white},
  rulecolor=\color{black}, numberstyle=\tiny}
\lstset{style=codigo}
\renewcommand{\lstlistingname}{Código}            % rótulo da legenda: "Código N"
\renewcommand{\lstlistlistingname}{Lista de códigos}
%%END:estilo
